% AI-produced by GPT-6 Astra (OpenAI, ChatGPT Pro mode), 2026-11-19.
% This segment uses only amsmath, amssymb, amsthm, and standard LaTeX commands.
% Bibliographic keys are supplied in B_references.bib.

\section{B.1. What the two-stratum sheaf description actually requires}

We write $\operatorname{Loc}(T)$ for the category of locally constant sheaves
of sets on a topological space $T$. A local system is not assumed to have
finite stalks. We keep this distinction because it affects the answer.

\paragraph{B1.1 (Proved: unconditional closed--open recollement).}
Let $i:Y\hookrightarrow X$ be a closed inclusion of topological spaces and
let $j:U=X\setminus Y\hookrightarrow X$ be its complementary open inclusion.
There is an equivalence
\[
 \operatorname{Sh}(X)\simeq
 \{(L,E,\alpha):L\in\operatorname{Sh}(U),\ E\in\operatorname{Sh}(Y),\
             \alpha:E\longrightarrow i^{-1}j_*L\}.
 \tag{B1.1}
\]
Morphisms on the right are pairs of sheaf morphisms commuting with $\alpha$.
Its restriction to the full subcategory $\mathcal F$ of sheaves locally
constant on $U$ and $Y$ is obtained by imposing
$L\in\operatorname{Loc}(U)$ and $E\in\operatorname{Loc}(Y)$. No normal
triviality, local contractibility, separation, or finiteness assumption is
needed for this statement.

\begin{proof}
Put $T=i^{-1}j_*L$ and form the pullback in $\operatorname{Sh}(X)$
\[
 F=j_*L\mathbin{\times}_{i_*T}i_*E,
 \tag{B1.2}
\]
where the first map is the adjunction unit and the second is $i_*\alpha$.
At $x\in U$ the two sheaves supported on $Y$ have terminal stalk, so
$F_x=L_x$. At $y\in Y$ the pullback on stalks is
$T_y\times_{T_y}E_y=E_y$. The resulting closed--open comparison is
$\alpha$. Conversely, adjunction supplies a map from every sheaf $F$ to
the pullback constructed from $(j^{-1}F,i^{-1}F)$; its map on every stalk
is an isomorphism. Sheaves on a topological space have enough points,
so this is an isomorphism. The pullback universal property proves the
assertion on morphisms. Empty $Y$ or $U$ gives the corresponding identity
equivalence. This is the explicit form of closed--open recollement; compare
\cite[Example A.8.4]{LurieHA}.
\end{proof}

\paragraph{B1.2 (Proved under an exact comparison hypothesis).}
Suppose in addition that $N\subset X$ is an open neighbourhood of $Y$ with
a retraction $r:N\to Y$, and write $N^*=N\setminus Y$ and
$p=r|_{N^*}:N^*\to Y$. For $L\in\operatorname{Loc}(U)$ there is a natural
germ comparison
\[
 \theta_L:p_*(L|_{N^*})\longrightarrow i^{-1}j_*L.
 \tag{B1.3}
\]
Indeed a section over $p^{-1}W$ gives a germ along the open subset
$W\subset Y$ by using the neighbourhood $r^{-1}W$.
If every $\theta_L$ is an isomorphism, then
\[
 \mathcal F\simeq
 \{(L,E,\varphi):L\in\operatorname{Loc}(U),\ E\in\operatorname{Loc}(Y),\
       \varphi:p^{-1}E\longrightarrow L|_{N^*}\}.
 \tag{B1.4}
\]
More generally, without assuming (B1.3) to be an isomorphism, the precise
necessary and sufficient condition for this functor of systems to be an
equivalence over the fixed pair $(L,E)$ is that
\[
 \operatorname{Hom}_Y(E,p_*(L|_{N^*}))\longrightarrow
 \operatorname{Hom}_Y(E,i^{-1}j_*L)
 \tag{B1.5}
\]
be bijective. Requiring the description after restriction to every open
subset of $Y$, and for every locally constant $E$, makes the isomorphism
condition on (B1.3) necessary as well as sufficient.

\begin{proof}
Adjunction identifies the set of maps $\varphi$ with the source of
(B1.5). Under (B1.1), its image is exactly the right-hand gluing map.
Thus bijectivity on each pair of objects is exactly full faithfulness
and essential surjectivity over the two restriction categories. It also
gives the assertion for varying objects, since morphisms are the pairs
commuting with their gluing maps. For the local necessity assertion we
take the constant one-point sheaf on each open $W\subset Y$: maps from
that sheaf are sections over $W$. Bijectivity for all $W$ is precisely
that $\theta_L$ is an isomorphism of sheaves.
\end{proof}

The usual two-stratum conical tube satisfies this criterion. More
explicitly, suppose $Y$ is paracompact and has a basis of contractible
open sets, that
$r:N\to Y$ is locally a bundle of open cones on compact, locally
contractible links, and that $U$ is locally path connected and semilocally
simply connected. We allow disconnected links. In a contractible base
chart $W$, the punctured tube has the form
$W\times K\times(0,\epsilon)$. A local system there is pulled back from
$K$, and its sections are unchanged by shrinking $W$ and $\epsilon$.
Compactness of $K$ makes the constant-radius conical neighbourhoods
cofinal at each point of $Y$. Consequently (B1.3) is an isomorphism on
stalks. Radial contraction also identifies
$\operatorname{Loc}(N)$ with $\operatorname{Loc}(Y)$. Monodromy now gives
the groupoid formulation in 151/035--036, with
$\Pi_1N\simeq\Pi_1Y$. Here monodromy means transport in an actual
locally constant sheaf; the assumptions on the spaces make this an
equivalence with all set-valued functors on their fundamental groupoids.

\paragraph{B1.3 (Refuted: an unconditional uniform pro-tube formula).}
The assertion on 151/035 is false if, without normal or neighbourhood
stability hypotheses, a map on a pro-tube means a morphism represented
on one sufficiently small neighbourhood. This fails even for a compact
metric $X$ with two strata, one a point and the other a disjoint union
of one-dimensional manifolds.

\begin{proof}
Let $X$ be the Hawaiian earring in $\mathbb R^2$: its $n$th circle $C_n$
has centre $(1/n,0)$ and radius $1/n$. Let $Y=\{o\}$, where $o=(0,0)$,
and $U=X\setminus\{o\}$. Each $C_n\setminus\{o\}$ is an open-and-closed
component of $U$, homeomorphic to an open interval. Let $L$ be the
constant two-point local system on $U$, and let $E=\mathbb N$ on $Y$.
For each $n$, choose a sufficiently small ball about $o$ that misses the
point of $C_n$ opposite $o$. On its punctured intersection with $X$,
there is a locally constant function that is $1$ on the upper arm of
$C_n$, $0$ on its lower arm, and $0$ on all other circles. Denote its
germ by $a_n\in(i^{-1}j_*L)_o$. The map $\alpha(n)=a_n$ defines an
object of (B1.1), hence a sheaf $F\in\mathcal F$.

Suppose it came from the manuscript's ordinary tube system. A locally
constant sheaf on the tube, with stalk $\mathbb N$ at $o$, is constant
on some single open neighbourhood $W$ of $o$. The map to $L$ on the
deleted tube would therefore be represented on a smaller such $W$.
Every neighbourhood of $o$ contains an entire $C_n$ for all sufficiently
large $n$. On the connected set $C_n\setminus\{o\}$, a morphism from
the constant $\mathbb N$-sheaf to the constant two-point sheaf sends
each fixed element to a constant value. This contradicts the two
distinct arm-values in $a_n$. Relabelling the stalk $\mathbb N$ does
not change the contradiction. No uniform representative exists.

There are no omitted hypotheses in this counterexample: $X$ is compact
metric; $Y$ is closed; $U$ is a disjoint union of open intervals; and
both restrictions of $F$ are locally constant by (B1.2). What fails is
uniform stabilization of neighbourhood gluing. In particular, the
unconditional statement (B1.1), not (B1.4), is the correct replacement.
The example also explains why one cannot assert that the pro-neighbourhood
of an arbitrary closed subset is homotopy equivalent to that subset.
\end{proof}

There is a second independent proviso in the wording of 151/035:
ordinary fundamental groupoids do not classify all local systems on
arbitrary spaces. For example, on a nondiscrete Cantor space the
fundamental groupoid has only identity arrows, so a representation may
assign a singleton to one non-isolated point and a two-point set to
every other point. No locally constant sheaf has those stalk
cardinalities. Thus the purely groupoid wording also requires a
monodromy hypothesis, even in the one-stratum case.

\paragraph{B1.4 (Proved: the inertia formula, with its precise scope).}
In the connected case let $D=\pi_1(N^*)$, $B=\pi_1(Y)$ and
$G=\pi_1(U)$, with homomorphisms $p_*:D\to B$ and $i_*:D\to G$.
Assume $p_*$ is onto and put $I=\ker p_*$. Then (18) of 151/036 is
equivalent to the $B$-equivariant map
\[
 E\longrightarrow (i_*^*L)^I.
 \tag{B1.6}
\]
Indeed a $D$-equivariant map from a $D$-set on which $I$ acts trivially
has image in the $I$-fixed points, and the induced action on those fixed
points factors through $D/I=B$. Conversely such a $B$-map is
$D$-equivariant. If $p$ is a fibration with connected fibre $K$, its
homotopy exact sequence gives
\[
 \pi_2(Y)\longrightarrow\pi_1(K)\longrightarrow D
      \longrightarrow B\longrightarrow1.
 \tag{B1.7}
\]
Thus $I$ is the quotient of $\pi_1(K)$ by the image of the connecting
homomorphism, and it equals $\pi_1(K)$ only when that homomorphism is
zero. This independently rechecks R4, SA10-C3-04. If $Y$ and $K$ are
aspherical, the same long exact sequence gives $\pi_q(N^*)=0$ for
$q\geq2$. Disconnected fibres are handled by the groupoid formulation;
surjectivity of $D\to B$ must not then be silently asserted.

\paragraph{B1.5 (Proved: the exact criterion for pro-tubes).}
We give a criterion that does not assume a chosen retraction or a
stationary representative. Let $\mathcal N_Y$ be the filtered ordered
set of open neighbourhoods of $Y$, with arrows given by shrinking.
Form the filtered bicategorical colimits in categories
\[
 \begin{aligned}
 \mathcal T&=\operatorname*{2colim}_{N\in\mathcal N_Y}
                         \operatorname{Loc}(N),\\
 \mathcal T^*&=\operatorname*{2colim}_{N\in\mathcal N_Y}
                         \operatorname{Loc}(N\setminus Y).
 \end{aligned}
 \tag{B1.8}
\]
Their objects and arrows have representatives on one neighbourhood;
two representatives agree if they agree after a common shrinking.
There are restriction functors
$r:\mathcal T\to\operatorname{Loc}(Y)$,
$a:\mathcal T\to\mathcal T^*$, and
$b:\operatorname{Loc}(U)\to\mathcal T^*$.
For $E\in\mathcal T$ and $L\in\operatorname{Loc}(U)$, restriction of
a representative morphism to germs along $Y$ defines
\[
 \mu_{E,L}:\operatorname{Hom}_{\mathcal T^*}(aE,bL)
 \longrightarrow
 \operatorname{Hom}_{Y}(rE,i^{-1}j_*L).
 \tag{B1.9}
\]
The natural functor from systems $(L,E,\varphi:aE\to bL)$ to
$\mathcal F$ is an equivalence, compatible with restriction to $U$
and $Y$, if and only if both of the following conditions hold:
$r$ is an equivalence, and every map (B1.9) is bijective.
These are respectively an extension condition for local systems and a
uniform extension condition for their gluing morphisms. For a
cofinal system of actual tube retractions, B1.2 supplies sufficient
geometric hypotheses for both conditions.

\begin{proof}
For a representative $E_N$ and a representative
$\varphi:E_N|_{N'\setminus Y}\to L|_{N'\setminus Y}$ with $N'\subset N$,
adjunction gives $E_N|_{N'}\to j'_*L|_{N'\setminus Y}$.
Its inverse image on $Y$ is (B1.9); further shrinking gives the same
map. This proves that $\mu$ and the asserted functor are well defined.
If $r$ is an equivalence, choose a quasi-inverse and normalize the
systems so that their second entry is an object of
$\operatorname{Loc}(Y)$. The two categories are then Grothendieck
constructions over the same category
$\operatorname{Loc}(U)\times\operatorname{Loc}(Y)$, with fibres the
sets on the two sides of (B1.9). Their arrow conditions are equality
of the composites along the two restrictions. The natural bijections
$\mu$ identify those conditions, proving sufficiency.

Conversely, set $L$ equal to the terminal local system on $U$.
There is a unique possible $\varphi$, so the systems with this
restriction form $\mathcal T$; the corresponding sheaves form
$\operatorname{Loc}(Y)$ by B1.1. Compatibility with restrictions
therefore forces $r$ to be an equivalence. After the preceding
normalization, the comparison has identical base objects and arrows.
Its fibre over a fixed $(L,E)$, with the identity base arrow, is a
discrete category. Both projections allow transport along every
isomorphism of base pairs: conjugate the gluing arrow by the two
given isomorphisms. Thus an equivalence over the base is an
equivalence on these fibres. This is exactly bijectivity
of $\mu_{E,L}$. This proves necessity.
\end{proof}

For groupoid wording, one must additionally impose the monodromy
equivalence on the spaces in a cofinal neighbourhood system and on
$U,Y$, or explicitly define the groupoids as those classifying these
local systems. Ordinary $\Pi_1$ is sufficient on locally path connected,
semilocally simply connected spaces. The pro-representation convention
in (B1.8) keeps the one-stage representability of arrows; it must not
be replaced silently by independent choices of a neighbourhood for
each element of an infinite stalk. In B1.3, $r$ is an equivalence
because every local system near the single point is constant after
shrinking. Precisely the surjectivity of $\mu_{\mathbb N,L}$ fails.

\section{B.2. Shape and the level at which gluing takes place}

\paragraph{B2.1 (Known, with an explicit derivation: the correct shape theorem).}
Let $\mathcal S$ be the $\infty$-category of spaces and
$\operatorname{Top}_\infty$ the $\infty$-category of $\infty$-topoi and
geometric morphisms. The shape $\Pi_\infty\mathcal X\in
\operatorname{Pro}(\mathcal S)$ is characterized by
\[
 \operatorname{Map}(\Pi_\infty\mathcal X,K)
  \simeq \Gamma_{\mathcal X}\Delta_{\mathcal X}K,
       \qquad K\in\mathcal S.
 \tag{B2.1}
\]
For every small diagram whose colimit in $\operatorname{Top}_\infty$
is $\mathcal X$, the canonical map
\[
 \operatorname*{colim}_a\Pi_\infty\mathcal X_a
       \longrightarrow\Pi_\infty\mathcal X
 \tag{B2.2}
\]
is an equivalence of pro-spaces. No stratification or normal
fibration hypothesis occurs here.

\begin{proof}
For a space $K$, the universal property of the slice $\infty$-topos
$\mathcal S_{/K}$ identifies
$\operatorname{Map}_{\operatorname{Top}_\infty}
(\mathcal X,\mathcal S_{/K})$ with the right side of (B2.1).
Mapping a colimit to this slice gives
\[
 \operatorname{Map}(\Pi_\infty\mathcal X,K)
 \simeq\lim_a\operatorname{Map}(\Pi_\infty\mathcal X_a,K).
\]
The fully faithful realization of pro-spaces as accessible left exact
functors $\mathcal S\to\mathcal S$, with variance reversed, detects
equivalences by these evaluations. This proves (B2.2). Equivalently,
shape is a left adjoint, as stated in
\cite[Remark 7.1.6.15]{LurieHTT} and
\cite[Definition 4.2.1]{BGHExodromy}.
\end{proof}

For an actual open cover of $X$, descent for sheaves of spaces gives
the colimit of sheaf $\infty$-topoi indexed by its full descent diagram.
Thus (B2.2) applies to every honest open-cover stage of a pro-tube
construction. If the pieces are locally of singular shape, their shapes
are represented by their singular homotopy types. This gives the
homotopy-colimit statement used in topology. Passing from one stage to
the pro-system is legitimate as a statement about the pro-system of
these colimits; it is not permission to commute an arbitrary infinite
homotopy colimit with a pro-limit. Finite homotopy colimits, indexed
by finite simplicial sets such as a span, can be computed levelwise
after a common cofiltered reindexing. A category with finitely many
objects and arrows need not have finite nerve, so that observation
does not justify truncating an untruncated higher descent diagram
to $\Delta_{\leq2}$. We retain the full stagewise \v Cech diagram.

More explicitly, let $A$ be the cofiltered category of synchronized
neighbourhood choices in A.1, with an arrow $\alpha'\to\alpha$ for
a shrinking. Write $C_q^\alpha$ for the disjoint union of all
$(q+1)$-fold intersections of its actual open cover. The precise
higher statement is the natural equivalence
\[
 \begin{aligned}
 &\left(\operatorname*{colim}_{[q]\in\Delta^{\mathrm{op}}}
       \Pi_\infty\operatorname{Sh}_\infty(C_q^\alpha)\right)_{\alpha\in A}
 \\
 &\hspace{3em}\simeq
 \operatorname{const}_A\bigl(\Pi_\infty\operatorname{Sh}_\infty(X)\bigr)
 \end{aligned}
 \tag{B2.2a}
\]
in $\operatorname{Fun}(A,\operatorname{Pro}(\mathcal S))$.
The comparison is induced by the cover augmentations, and is an
equivalence at every $\alpha$, which proves the displayed natural
equivalence. If every shape in this formula is represented by the
corresponding singular homotopy type, it also says that the resulting
cofiltered system of homotopy colimits is the constant pro-space
$\operatorname{Sing}(X)$. No exchange of the two indexing operations
is made.

The Artin--Mazur qualification also matters. The coherent hypercover
construction models the shape of the \emph{hypercomplete} sheaf
$\infty$-topos: see \cite[Proposition 5.1]{HoyoisGalois}. It yields
the classical pro-homotopy object after the stated loss of coherence.
For the locally contractible spaces used in the examples below these
distinctions do not alter the answer. They cannot be erased for an
arbitrary topos; compare \cite[Example 7.1.6.9]{LurieHTT}.

\paragraph{B2.2 (Refuted: arbitrary bicategorical colimits of ordinary topoi).}
The same sentence is false if ``colimit of topoi'' means a bicategorical
colimit of ordinary Grothendieck topoi, and the right side means a
homotopy colimit of their shapes.

\begin{proof}
Let $G=C_2$ and let $BG$ denote the ordinary topos of $G$-sets. Its
unique geometric morphism $q:BG\to\mathbf{Set}$ has fully faithful
inverse image $q^*$, the trivial-action functor. The bicategorical
pushout of
\[
 \mathbf{Set}\longleftarrow BG\longrightarrow\mathbf{Set}
 \tag{B2.3}
\]
is $\mathbf{Set}$. To verify the universal property, fix any target
topos $\mathcal E$. A pseudocone consists of two geometric morphisms
$a,b:\mathbf{Set}\to\mathcal E$ and an invertible transformation
$q^*a^*\cong q^*b^*$. Full faithfulness of $q^*$ identifies that
transformation, and all transformations of pseudocones, with an
invertible transformation $a^*\cong b^*$ in $\mathbf{Set}$. The category
of two objects with a specified isomorphism between them is equivalent
to the original category by evaluation on the first object. This is
exactly the bicategorical pushout universal property.

The associated sheaf $\infty$-topos of $BG$ is
$\operatorname{Fun}(BG,\mathcal S)$, whose shape is the classifying
space of $G$; see \cite[Example 4.2.3]{BGHExodromy}, or compute global
sections of a constant diagram directly. Hence the homotopy pushout of
the three shapes is $\Sigma BC_2$. Its reduced second integral homology
is $H_1(BC_2;\mathbb Z)=\mathbb Z/2$, while the shape of the
bicategorical pushout is a point. The program
\texttt{check\_B\_finite.py} verifies the normalized bar differentials
$d_1=0,d_2=2,d_3=0,d_4=2$ and the finite group operations.
\end{proof}

This counterexample identifies the exact failure: passage from an
ordinary topos to its associated sheaf $\infty$-topos does not preserve
all bicategorical colimits. The actual open-cover gluings satisfy
descent at every higher level, so they are not affected by this
counterexample. We must prove that higher descent for the diagram
under consideration before applying B2.1.

\paragraph{B2.3 (Known under precise hypotheses: Quinn's input and L2).}
For a metric space $X$ and a closed subset $Y$, put
\[
 H=\operatorname{holink}(X,Y)
  =\{\gamma:[0,1]\to X:\gamma(0)\in Y,
                 \ \gamma(t)\notin Y\ (t>0)\}.
\]
The maps $e_0:H\to Y$ and $e_1:H\to U$ are endpoint evaluations.
If $Y$ is forward tame in $X$, then the ordinary homotopy types satisfy
\[
 X\simeq Y\mathbin{\cup^h}_{H}U.
 \tag{B2.4}
\]
Quinn proves a controlled strict relative version in
\cite[Lemma 2.4(2)]{Quinn1988}, using his controlled holink and
mapping-cylinder convention. We give a separate proof of (B2.4)
that does not identify that cylinder topology with the ordinary
quotient-cylinder topology. Let $N$ range over the open neighborhoods
of $Y$. The numerable open cover $X=N\cup U$ gives, naturally in $N$,
the homotopy-pushout equivalence
$N\cup^h_{N\setminus Y}U\simeq X$.
A fixed tame deformation gives an equivalence
$\{N\}_N\simeq\operatorname{const}Y$ in prospaces: for any target
neighborhood $N$, shrink to the points whose entire deformation lies
in $N$. The same fixed deformation and the explicit endpoint-preserving
path homotopy of C4.02 give
$\{N\setminus Y\}_N\simeq\operatorname{const}H$.
Those formulas need only tameness of this pair; the fibration part
of the HSS definition is not used. The comparison uses the natural
zigzag through paths wholly contained in $N$, so its endpoint maps
and the homotopy between the two maps to $X$ are the original ones.
The finite span colimit in $\operatorname{Pro}(\mathcal S)$ is
computed on these synchronized levels, because filtered colimits
of spaces commute with finite simplicial limits. Replacing the two
pro-germs by their constant equivalents therefore gives
$\operatorname{const}X\simeq
\operatorname{const}(Y\cup^h_HU)$.
The constant embedding is fully faithful, proving (B2.4).
In the usual quotient double-cylinder model, the associated map to
$X$ is induced by the continuous path evaluation
$H\times[0,1]\to X$.
Quinn's Definition 3.1 adds, for each incident pair of strata, that
$e_0$ is a fibration. That is the precise normal-fibration analogue of
151/032, once a deleted tube has been compared with the holink. It is
not the whole definition: forward tameness remains a separate
hypothesis. The fibration condition alone does not imply tameness.
For the Hawaiian earring with $Y$ its common point, $e_0$ is a fibration
because its target is a point. But $Y$ is not tame: a neighbourhood
deformation into $Y$ would null-homotope every sufficiently small
circle in $X$, while retraction of $X$ onto any one of those circles
shows its fundamental loop is nonzero. Consequently L2 must be read
with this additional tameness and comparison data, and its unrestricted
topos-colimit conclusion must be replaced by B2.1. The shape theorem
itself needs neither condition.

\paragraph{B2.4 (Proved: rechecking the two examples in R5).}
For $(\mathbb C,\{0\})$, B1.4 identifies a constructible sheaf with a
set $E$, a $\mathbb Z$-set $M$, and a map $E\to M^{\mathbb Z}$.
Equivalently it is a covariant set-valued functor on the category with
objects $p,q$, automorphism group $\mathbb Z$ at $q$, and a unique map
$p\to q$. The object $p$ is initial; hence the classifying space of
this category, and the shape of its functor topos, is contractible.
This agrees with $\mathbb C$.

For $(S^2,\{p\})$, the complementary plane is contractible. The same
description reduces to a single map of sets $E\to M$. The corresponding
arrow category also has contractible classifying space. Nevertheless
the genuine tube span has shapes
\[
 *\longleftarrow S^1\longrightarrow *,
\]
whose homotopy pushout is $\Sigma S^1\simeq S^2$. Thus the ordinary
constructible-sheaf topos does not retain the sphere's higher homotopy
type. This proves the assertions SA10-C3-05 and SA10-C3-06 without
identifying that functor topos with the required higher gluing.

\section{B.3. Exit paths, truncation, and gerbe gluing}

\paragraph{B3.1 (Known; truncation proved below).}
Assume that $X$ is paracompact, locally of singular shape, and conically
stratified by a poset $P$ satisfying the ascending chain condition.
These hypotheses hold, for example, for a finite-dimensional,
paracompact Hausdorff topologically stratified space with the compact
conical charts used in \cite[Definition 3.1]{Treumann2009}.
Write $\operatorname{Exit}(X)=\operatorname{Sing}_P(X)$.
For every integer $n\geq0$ there are equivalences
\[
 \begin{split}
 \operatorname{Cons}_{\leq n}(X)
 &\simeq \operatorname{Fun}(\operatorname{Exit}(X),\mathcal S_{\leq n})\\
 &\simeq \operatorname{Fun}(h_{n+1}\operatorname{Exit}(X),
                                      \mathcal S_{\leq n}).
 \end{split}
 \tag{B3.1}
\]
Here the left side consists of constructible $n$-truncated sheaves of
spaces, and $h_{n+1}$ retains the objects and the $n$-truncations of
all mapping spaces, with their coherent compositions.
The regularity hypotheses on the strata are substantive. If
``conical'' only means a local product with a cone, allowing an
arbitrary unstratified base, every one-stratum space is conical.
The Cantor-space example after B1.3 then refutes the unqualified
claim already at $n=0$, and also at every higher level by regarding
sets as discrete spaces. For the usual finite Thom--Mather,
Whitney, or conically smooth spaces with manifold strata, the
hypotheses used here hold. Disconnected and empty strata cause no
change: the statement is componentwise and an empty stratum
contributes no objects.

\begin{proof}
The untruncated equivalence is
\cite[Theorem A.9.3]{LurieHA}. Its functor evaluates on stalks as the
corresponding functor evaluates on objects, by
\cite[Proposition A.9.16]{LurieHA}. A constructible object is
$n$-truncated precisely when all these stalks are $n$-truncated:
constructible objects here are hypercomplete, and the assertion also
follows directly from their local conical descriptions. Thus the
equivalence restricts to the first line of (B3.1); this is also recorded
in \cite[Remark A.9.7]{LurieHA}. Mapping spaces in
$\mathcal S_{\leq n}$ are $n$-truncated, since they are limits of
$n$-truncated spaces. The defining universal property of the homotopy
$(n+1)$-category therefore gives the second equivalence, including its
mapping spaces and not merely its objects.
\end{proof}

We give the finite-stratum induction explicitly, including the descent
reduction; this avoids identifying a global undercategory with a single
local link. We prove the untruncated comparison and simultaneously all
its truncated restrictions by induction on the number of nonempty
stratum labels. The empty case is terminal, and the one-stratum case
is ordinary monodromy, \cite[Theorem A.4.19]{LurieHA}.
For an open $F_\sigma$ subset $W\subset X$, let $\mathcal B_W$ be the
category of its conical open $F_\sigma$ charts, as in
\cite[Lemma A.5.16]{LurieHA}. The opposite of the category of such charts
containing any fixed point is filtered, since smaller conical charts
refine intersections. Thus its nerve is contractible. Exit excision
\cite[Theorem A.7.1]{LurieHA} and descent for constructible sheaves give
\[
 \begin{split}
 \operatorname{Fun}(\operatorname{Exit}(W),\mathcal S)
 &\simeq\lim_{V\in\mathcal B_W^{\mathrm{op}}}
              \operatorname{Fun}(\operatorname{Exit}(V),\mathcal S),\\
 \operatorname{Cons}(W)
 &\simeq\lim_{V\in\mathcal B_W^{\mathrm{op}}}\operatorname{Cons}(V).
 \end{split}
 \tag{B3.2}
\]
The same holds for $n$-truncated values, since truncation is local.
The comparison is natural for these restrictions, so an equivalence
on the charts implies an equivalence on $W$ by (B3.2).

It remains to prove the induction step on a chart
$V=Z\times C(K)$, with closed stratum $Z$ and
$V^*=Z\times K\times(0,\infty)$. Here $V^*$ has fewer nonempty
stratum labels. We spell out the categorical induction criterion, so
that no global nearby-cycle stability is presumed.
In $\mathcal C=\operatorname{Fun}(\operatorname{Exit}(V),\mathcal S)$
let $\mathcal C_0$ be functors extended by the initial space from
$\operatorname{Exit}(V^*)$, and let $\mathcal C_1$ be functors sending
every arrow to an equivalence. In
$\mathcal D=\operatorname{Cons}(V)$ let $\mathcal D_0=j_!
\operatorname{Cons}(V^*)$ and let $\mathcal D_1$ be the local systems
pulled back from $Z$. The inclusions of each of these pairs have
complementary right adjoints. On the functor side this follows from
the cone contraction and the dual of \cite[Lemma A.9.18]{LurieHA};
on the sheaf side it is \cite[Lemma A.9.13]{LurieHA}.
The dual of the recollement comparison criterion
\cite[Proposition A.8.14]{LurieHA} now applies for the following
explicit reasons. The restriction
$\mathcal C_0\to\mathcal D_0$ is the comparison on $V^*$, hence an
equivalence by induction. The restriction
$\mathcal C_1\to\mathcal D_1$ is monodromy on $Z$, because radial
contraction identifies the locally constant objects on $V$ with
those on $Z$; the precise sheaf comparison is
\cite[Lemmas A.9.11--A.9.12]{LurieHA}. The comparison preserves
pushouts by \cite[Lemma A.9.15]{LurieHA}. Finally it carries the
$\mathcal C_0$-colocalization of an object of $\mathcal C_1$ to the
corresponding $\mathcal D_0$-colocalization: both are obtained by
restriction to $V^*$ followed by extension by the initial object.
This verifies every hypothesis of the criterion, proving
$\mathcal C\simeq\mathcal D$ and completing the induction.

All these comparisons respect stalks and evaluation, so the full
subcategories whose values are $n$-truncated agree at each induction
step. Mapping into $\mathcal S_{\leq n}$ then factors uniquely through
$h_{n+1}\operatorname{Exit}(X)$ by the already proved universal
property. This proves every truncation level by the same induction.
For infinite ACC posets, the source theorem replaces the finite
induction by rank induction; finite induction alone would not prove
that extension.

The associated closed--open gluing operation is
\[
 (F_U,F_S,\alpha:F_S\longrightarrow i^*j_*F_U).
 \tag{B3.3}
\]
It is the homotopy-pullback version of (B1.2). In an actual conical
chart with contractible base and compact link $K$, its nearby value
is $\lim_{\operatorname{Exit}(K)}F_U$. On the exit side this is the
right Kan extension from the open subcategory. The open cone adds
an initial cone point. For conically smooth charts the exact cone
and sheaf comparisons are \cite[Lemmas 6.1.4 and 6.1.7]{AFTLocal},
taking opposites by its Corollary 1.2.10 to convert enter paths to exit paths;
the general topological conical case is included in the foregoing
recollement argument. Limits of $n$-truncated spaces remain
$n$-truncated, which also makes the compatibility of (B3.3) with
every finite level explicit.

\paragraph{B3.2 (Proved: groupoid stacks and the correction to L3).}
At $n=1$, (B3.1) is the equivalence between constructible stacks in
groupoids and homotopy-coherent $2$-functors from the exit-path
$2$-category to groupoids. It is also the restriction of
\cite[Theorem 7.14]{Treumann2009} from category-valued to
groupoid-valued stacks: the value at each object is the stalk
category, so the condition that these values are groupoids is
preserved in both directions. His Theorem 7.13 proves compatibility
with open-cover descent.

The set-valued level, by contrast, is $n=0$: it retains $\pi_0$ of
each exit mapping space. Thus L3's phrase ``truncated at $\pi_1$ of
the links'' is one level too high if it means that link $\pi_1$
survives as automorphisms of exiting arrows. Those automorphisms
first appear in the stack level. For $(S^2,\{p\})$, the two-stratum
exit category can be represented by two objects and cross mapping
space $S^1$: the open stratum is contractible, and its attaching
holink has the homotopy type of $S^1$. Taking $\pi_0$ makes a single
arrow, while the fundamental groupoid of that mapping space retains
$\mathbb Z$. This is exactly the difference seen by the gerbes.

One should also distinguish two uses of gluing. A general constructible
stack is a closed--open system (B3.2), with a possibly noninvertible
functor to nearby data. An $A$-banded gerbe is locally equivalent to
$BA$ at every point of $X$. Its gluing over an honest open tube and
$U$ uses an \emph{equivalence} on their overlap, as in 151/042.
These gerbes form a subproblem of the constructible-stack theorem;
arbitrary constructible stacks are not all gerbes.

\paragraph{B3.3 (Proved: the gerbe computation, including the quotient by gauges).}
Let $A$ be any discrete abelian group with constant band. Then
\[
 H^2(S^2;A)\cong A.
 \tag{B3.4}
\]
We obtain this by the precise gluing on 151/041--043.

\begin{proof}
Cover $S^2$ by $U=S^2\setminus\{p\}$ and an open disk $N$ about $p$.
Both are contractible, and $U\cap N$ is an annulus. Every $A$-banded
gerbe on either contractible piece is equivalent to the neutral
gerbe $BA$. An equivalence between two copies of $BA$ on an overlap
is an $A$-torsor there; composition is contracted product of torsors.
The isomorphism classes of these torsors are
\[
 H^1(U\cap N;A)=\operatorname{Hom}(\mathbb Z,A)=A.
\]
Changing either neutralization acts by the restrictions of torsors
on $U$ and $N$. Those torsors have trivial isomorphism class because
$H^1(U;A)=H^1(N;A)=0$. Consequently this action identifies no two
distinct elements of $A$. Open-cover $2$-descent is effective, so
every such torsor glues the two neutral gerbes, and two glued gerbes
are equivalent precisely when the torsor classes agree. This proves
(B3.4), using the interpretation of $H^2$ by banded gerbes; the
precise abelian classification used here is
\cite[Theorem 2.7]{DeopurkarGerbes}. The identification of
band-preserving self-equivalences of the neutral gerbe with torsors
is also explicit in its Remark 3.1: the image of the trivial torsor
determines the equivalence, which is contracted product with that
torsor (or its inverse according to the chosen action convention).

In the extension language of the manuscript, both pieces yield
trivial central extensions. An automorphism of the pulled-back
extension $A\times\mathbb Z\to\mathbb Z$, fixing the band and the
quotient, is uniquely $(a,m)\mapsto(a+mh,m)$ for $h\in A$.
Thus the same gluing set is $\operatorname{Hom}(\mathbb Z,A)$.
Choices of basepoints and neutralizations are why this is a
description up to equivalence of gerbes, not a canonical literal
identity between a gerbe and a group extension. This rechecks
SA10-C3-07.
\end{proof}

\section{B.4. The aspherical programme and the coherence it needs}

\paragraph{B4.1 (Proved: the premise in the question is refuted).}
For a two-stratum homotopy-pushout decomposition
$X\simeq Y\cup^h_H U$, suppose every component of $Y,H,U$ is a
$K(\pi,1)$ of CW homotopy type. Then the diagram
\[
 \Pi_1Y\longleftarrow\Pi_1H\longrightarrow\Pi_1U
 \tag{B4.1}
\]
\emph{does} determine the full weak homotopy type of $X$:
\[
 X\simeq B\Pi_1Y\mathbin{\cup^h}_{B\Pi_1H}B\Pi_1U.
 \tag{B4.2}
\]
No extra invariant is required for this span. Thus R5 does not
establish the insufficiency asserted in the last sentence of B.4.

\begin{proof}
For every space $T$ there is a natural map of simplicial sets
$\operatorname{Sing}(T)\to N\Pi_1(T)$ to its fundamental groupoid
nerve. If every component of $T$ is aspherical, this map is a weak
equivalence. Apply it objectwise to the span. A natural objectwise
equivalence of diagrams induces an equivalence on homotopy
colimits; composing with the hypothesized homotopy-pushout
comparison gives (B4.2). This is a natural zigzag of weak
equivalences, so it does not require arbitrary basepoint identifications.

For the sphere, (B4.1) is $1\leftarrow\mathbb Z\to1$, and (B4.2)
is $*\cup^h_{S^1}*\simeq S^2$. It is the set-valued functor topos
that loses this information. On 151/041 Grothendieck explicitly
distinguishes these assertions and withdraws the feared failure of
groupoid-span reconstruction. The prompt's interpretation, not
that sentence of the manuscript, requires correction.
\end{proof}

The hypothesis on $H$ follows from the programme's local hypotheses
when $e_0:H\to Y$ is a fibration with aspherical fibres and $Y$ is
aspherical: the homotopy exact sequence proves asphericity of each
component of $H$. This works with disconnected fibres as well;
groupoids avoid an unjustified surjectivity assertion on $\pi_1$.

\paragraph{B4.2 (Proved: the full groupoid assembly data).}
Let $X\to P$ satisfy the hypotheses of B3.1, and put
$\mathcal E=\operatorname{Exit}(X)$. For every nonempty strict chain
$\sigma=(p_0<\cdots<p_m)$ in $P$, define
\[
 D(\sigma)=\operatorname{Map}_{/P}(\sigma,\mathcal E).
 \tag{B4.3}
\]
This is a space: a natural transformation over the fixed chain has
components inside strata, hence all its components are equivalences.
Restriction along subchains defines a homotopy-coherent functor
$D:\operatorname{sd}(P)^{\mathrm{op}}\to\mathcal S$.
Without any asphericity assumption, the complete assembly formula is
\[
 \operatorname{Sing}(X)\simeq
 \operatorname*{colim}_{\sigma\in\operatorname{sd}(P)^{\mathrm{op}}}
 D(\sigma).
 \tag{B4.3a}
\]
Assume that every stratum $D(p)$ and every total incidence link
$D(p<q)$ is a $1$-type. Then the groupoids
\[
 G_\sigma=\Pi_1D(\sigma),
 \qquad
 G:\operatorname{sd}(P)^{\mathrm{op}}\longrightarrow\operatorname{Gpd},
 \tag{B4.4}
\]
with \emph{all} restriction functors and their coherent comparison
isomorphisms determine the full homotopy type:
\[
 \operatorname{Sing}(X)\simeq
       \operatorname*{colim}_{\sigma\in\operatorname{sd}(P)^{\mathrm{op}}}
                  BG_\sigma.
 \tag{B4.5}
\]
The colimit in (B4.5) is taken in spaces, not in ordinary groupoids.
The incidence-link hypothesis holds when the local homotopy links
are aspherical and their source maps are fibrations over
aspherical strata. In geometric applications, ``all links'' must
mean all incident pairs, not merely adjacent pairs or a union of
different link strata.

\begin{proof}
The Segal equivalence gives
\[
 D(p_0<\cdots<p_m)\simeq
 D(p_0<p_1)\mathbin{\times^h}_{D(p_1)}\cdots
       \mathbin{\times^h}_{D(p_{m-1})}D(p_{m-1}<p_m).
 \tag{B4.6}
\]
This follows by mapping the concatenation of the elementary edges,
which is the chain $\sigma$ in $\operatorname{Cat}_\infty$, into
$\mathcal E$ over $P$; it is
\cite[Lemma 2.7.3]{BGHExodromy}. Homotopy limits of $1$-types are
$1$-types. Thus every value of $D$ is a $1$-type and the natural
objectwise truncation $D\to BG$ is an equivalence of diagrams.

To identify its colimit, extend $D$ along strict chains to all
simplices of $P$. Repeating an index inserts equivalences inside a
stratum, and this extension is the complete Segal nerve of
$\mathcal E$; this is the construction and proof of
\cite[Theorem 2.7.4]{BGHExodromy}. Its geometric realization is
the groupoid completion $|\mathcal E|$. Colimits commute with this
left Kan extension, so
$\operatorname{colim}_{\operatorname{sd}(P)^{\mathrm{op}}}D
\simeq|\mathcal E|$. Finally the natural inclusion
$\operatorname{Exit}(X)\to\operatorname{Sing}(X)$ induces an
equivalence on groupoid completions by
\cite[Corollary A.9.4]{LurieHA}. These three specified comparisons
give (B4.5).

For conical or Quinn pairs, total incidence links are modeled by
the endpoint-variable homotopy-link spaces. The latter's source
fibration has the local normal link as fibre. For popath models
the comparison that fixes both endpoints is
\cite[Theorem 4.9]{Miller2009}. The long exact sequence of that
fibration proves the last assertion of the theorem.
\end{proof}

The proof of (B4.3a) is the colimit and groupoid-completion part of
the same argument, before truncating $D$. This also specifies the
answer if ``links are $K(\pi,1)$'' refers only to the whole normal
link of each stratum. That condition does not imply our pairwise
incidence hypothesis. For example, stratify the closed $4$-ball by
a boundary point $p$, the remaining boundary $S^3\setminus\{p\}$,
and its interior. All three strata are contractible; the whole
normal link at $p$ is a $3$-ball and the normal link along the
remaining boundary is a point. But the incidence link from $p$
into the boundary stratum is $S^2$, not a $1$-type. In this reading,
one retains the spaces $D(\sigma)$, equivalently their full
Postnikov towers with the bonding maps, $k$-invariants and coherent
restriction maps. Formula (B4.3a) proves that this is sufficient in
every case. Formula (B4.6) says that the higher information lies in
the incidence links and their composition maps, not in independent
new spaces for each longer flag. Under the stronger pairwise
asphericity reading, all those towers truncate to the explicit
groupoid data (B4.4).

Formula (B4.4) is the requested additional data in the many-stratum
case. Here $\operatorname{Gpd}$ means its $(2,1)$-category, or
equivalently $\mathcal S_{\leq1}$, and not the category obtained by
discarding natural isomorphisms. Explicitly, the datum consists of
groupoids $G_\sigma$, functors
$F_{\sigma\tau}:G_\tau\to G_\sigma$ for inclusions
$\sigma\subset\tau$, and natural isomorphisms
$\theta_{\sigma\tau\upsilon}:F_{\sigma\tau}F_{\tau\upsilon}
\Rightarrow F_{\sigma\upsilon}$. Identity restrictions are identities,
the corresponding comparisons are identities, and for
$\sigma\subset\tau\subset\upsilon\subset\nu$ we require the equality
\[
 \theta_{\sigma\tau\nu}\circ
          (F_{\sigma\tau}\theta_{\tau\upsilon\nu})
 =\theta_{\sigma\upsilon\nu}\circ
          (\theta_{\sigma\tau\upsilon}F_{\upsilon\nu}).
 \tag{B4.8}
\]
Both sides have source
$F_{\sigma\tau}F_{\tau\upsilon}F_{\upsilon\nu}$ and target
$F_{\sigma\nu}$. These are all the coherence conditions, since
natural transformations between functors of groupoids form sets;
there is no additional layer of higher modifications.
An equivalent Segal presentation is: stratum and incidence groupoids;
composition on triple flags; the natural isomorphisms comparing
successive restrictions; and their coherence on longer flags.
Taking one object per connected groupoid converts functors into
homomorphisms only after choices, and converts the comparison
isomorphisms into conjugating elements. Those elements must be kept.
A diagram merely in the homotopy category of groups or spaces
forgets them. Conversely the full coherent diagram above is a
complete answer because (B4.5) reconstructs $X$, and it is recovered
from $X$ by (B4.3)--(B4.4).

\paragraph{B4.3 (Proved: why coherence cannot simply be discarded).}
The trivial circle bundle $S^2\times S^1\to S^2$ and the Hopf
bundle $S^3\to S^2$ have locally the same aspherical pieces and the
same maps on fundamental groups over a good cover. Their gluing
coherences differ, and their total spaces have different fundamental
groups.

\begin{proof}
Take a good cover of $S^2$ with nerve the boundary of a tetrahedron.
For example, use the open stars of the four vertices in its
geometric triangulation, with a small compatible thickening if
desired. Every nonempty intersection is contractible. Restrict
either oriented circle bundle to every such intersection: the
restriction is trivial and has homotopy type $S^1$. Every inclusion
between these pieces induces the identity homomorphism of
$\mathbb Z$, after orienting the fibre. Therefore the diagram
of individual groups and homomorphisms in the homotopy category
is the constant $\mathbb Z$-diagram for both bundles.

Choose transition maps $g_{ij}:U_i\cap U_j\to S^1$ and real lifts
$\theta_{ij}$ on the contractible pair intersections. On a triple
intersection,
\[
 n_{ijk}=\theta_{ij}+\theta_{jk}-\theta_{ik}\in\mathbb Z
 \tag{B4.7}
\]
is constant. These integers are the natural-transformation
coherence of the groupoid diagram; changing lifts adds a
coboundary. For the trivial bundle the class is zero. For the
Hopf bundle it is a generator: its north--south trivializations
have equatorial transition map of degree one, which is exactly
the winding number measured by (B4.7). This can also be seen
directly from the standard map $(z_0,z_1)\mapsto[z_0:z_1]$
on $S^3\subset\mathbb C^2$, whose two fibre coordinates differ
by the phase of $z_1/z_0$. The program
\texttt{check\_B\_finite.py} checks a tetrahedral representative:
the cocycle $(1,0,0,0)$ pairs to $-1$ with the fundamental cycle
$(-1,1,-1,1)$, so it is not a coboundary.

The total spaces have $\pi_1(S^3)=0$ and
$\pi_1(S^2\times S^1)=\mathbb Z$, respectively. Thus the
discarded coherence changes even the fundamental group of the
assembled space. To view this as a stratified example, lift the
finite face stratification of the tetrahedral sphere to either
bundle. Each stratum is $S^1$ times an open simplex and is
aspherical. Each local normal link is a point, a finite discrete
space, or a circle, hence componentwise aspherical. The same
coherence issue remains in passing from the full flag groupoid
diagram to its homotopy-category shadow.
\end{proof}

We do not claim that the full coherent span in B4.1 fails in this
example: a coarser two-open decomposition records the clutching
integer in its maps of fundamental groups. The loss occurs
precisely when the full coherent diagram is replaced by the list
of its local groups and homomorphisms up to homotopy.
